\section{構文と判断}

\subsection{型とスキーム}

capability $\kappa$ と target type $\tau$ は別 sort である．matcher と slot は同じ二つの
index を持つが，前者は producer，後者は consumer expectation を表す．
\[
\begin{array}{rcll}
\kappa &::=& \Any\mid\chi\mid\hat{s}
  \mid K\,\overline\kappa\mid(\kappa_1,\ldots,\kappa_n)
  & \kappa:\CapSort,\\
\tau &::=& \alpha\mid\hat a\mid\mathbf 1\mid\mathsf{Int}
  \mid\mathsf{Bool}\mid K\,\overline\tau
  \mid(\tau_1,\ldots,\tau_n)\mid\tau_1\to\tau_2
  & \\
&& \mid\Matcher{\kappa}{\tau}\mid\MatcherSlot{\kappa}{\tau}
  & \tau:\Ty,\\
\sigma &::=& \forall(\overline\chi:\CapSort).
  \forall(\overline\alpha:\Ty).\tau
  & \text{expression scheme},\\
\eta &::=& \kappa\mathbin{\triangleright}\tau
  & \text{pattern dual},\\
\varsigma &::=& \forall(\overline\chi:\CapSort).
  \forall(\overline\alpha:\Ty).
  (\overline\eta\Rightarrow\eta)
  & \text{pattern-function scheme}.
\end{array}
\]
$\chi,\alpha$ は flexible metavariable，$\hat s,\hat a$ は rigid skolem である．
$\Any$ は one-way capability matching の consumer 側でだけ wildcard として働き，
対称 unification では通常の rigid constructor である．$\mono\tau$ は量化 binder を持たない
scheme を表す．

\subsection{ソース構文}

\[
\begin{array}{rcl}
e &::=& x\mid\lambda x.e\mid e_1\,e_2
  \mid\mathtt{let}\ x=e_1\ \mathtt{in}\ e_2
  \mid\mathtt{fix}\ f\ x.e\mid n\\
&&\mid(e_1,\ldots,e_n)\mid C\,\overline e
  \mid op(\overline e)\mid\mathtt{something}\\
&&\mid\mathtt{matcher}\;\mathit{cls}\\
&&\mid\mathtt{matchAll}\ e_t\ \mathtt{as}\ e_m\
       \mathtt{with}\ p\to e,\\[0.2em]
p &::=& \$x\mid\_\mid\#e\mid\widetilde x
  \mid c\,\overline p\mid p_1\mathbin{\&}p_2
  \mid p_1\mathbin{|}p_2\mid f\,\overline p
  \mid(\overline p),\\
pp &::=& \$\mid\_\mid\#\$x\mid c\,\overline{pp}
  \mid(\overline{pp}),\\
dp &::=& \$x\mid\_\mid C\,\overline{dp}\mid(\overline{dp}),\\
cl &::=& pp\ \mathtt{as}\ M\ \mathtt{with}\
  [dp_j\to N_j]_{j=1}^{r},\\
\mathit{cls} &::=& [cl_i]_{i=1}^{m}.
\end{array}
\]
$c$ は pattern constructor，$C$ は data constructor である．surface
$\mathtt{match}$ は $\mathtt{matchAll}$ と $\mathsf{head}$ への elaboration であり，
core primitive ではない．pattern-function declaration は source expression の typing より前に
freeze され，canonical signature の一部になる．

\subsection{シグネチャ・文脈・状態}

\[
\begin{array}{c@{\qquad}l}
\Sighat=(\Sighat_D,\Sighat_P,\Sighat_F,
  \Sighat_{\mathit{prim}},\mathit{Obs}_{\Sighat})
  & \text{freeze 済み canonical signature},\\
\Gamma::=\varnothing\mid\Gamma,x:\sigma
  & \text{expression scheme context},\\
\Phi::=\varnothing\mid\Phi,x:\eta
  & \text{pattern-parameter context},\\
\Delta::=\varnothing\mid\Delta,x:\tau
  & \text{monomorphic pattern bindings},\\
q=(q_\kappa,q_\tau)
  & \text{next capability／target variable numbers},\\
S=(C,T)
  & \text{prevailing paired substitution}.
\end{array}
\]
\newpage
$q_\kappa$ と $q_\tau$ は独立な自然数 counter であり，それぞれ次に未使用の
capability metavariable と target metavariable の番号を保持する．以下では
\[
\chi_{q_\kappa}:\CapSort,
\qquad
\alpha_{q_\tau}:\Ty,
\qquad
q+\langle i,j\rangle=(q_\kappa+i,q_\tau+j)
\]
と書く．従って capability variable を一つ生成すれば supply は
$q+\langle1,0\rangle$，target variable なら $q+\langle0,1\rangle$，両方を一つずつ
生成する pattern variable 規則では $q+\langle1,1\rangle$ となる．番号空間は sort ごとに
別なので，$\chi_3$ と $\alpha_3$ は別の変数である．closed wrapper の初期 supply は，
signature と context に既出の番号より各 counter が大きくなるように選ぶ．

$S$ の作用は capability を先に，target substitution を後に適用する．
\[
S\kappa=C\kappa,
\qquad
S\tau=T(C\tau),
\qquad
S\Gamma=\{x:S\sigma\mid x:\sigma\in\Gamma\}.
\]
substitution delta は $\mathsf{seq}$ で時系列順に prevailing substitution へ合成する．
後段の capability action は前段の target range 内部にも作用するため，二 component を
独立に合成しない．

\subsection{判断の一覧}

\begin{center}
\small
\begin{tabular}{>{$}l<{$} l}
\toprule
\DDSynth_{\Sighat}(q,S,\Gamma,e,\tau,q',S')
  & expression synthesis \\
\DDCheck_{\Sighat}(q,S,\Gamma,e,\rho,q',S')
  & expression checking \\
\DDSynths_{\Sighat},\;\DDChecks_{\Sighat}
  & expression-list traversal \\
\DDPattern_{\Sighat}(q,S,\Gamma,\Phi,\Delta,p,
  \eta,\Delta',q',S')
  & user-pattern synthesis \\
\DDPPat_{\Sighat}(q,S,pp,\tau,\overline\eta,
  \Delta,q',S')
  & primitive-pattern checking \\
\DDDPat_{\Sighat}(q,S,dp,\tau,\Delta,q',S')
  & data-pattern checking \\
\DDClause_{\Sighat}(q,S,\Gamma,cl,\tau,
  \overline\eta,q',S')
  & matcher-clause inference \\
\DDArms_{\Sighat},\;\DDClauses_{\Sighat}
  & arm／clause-list traversal \\
\bottomrule
\end{tabular}
\end{center}
list judgment は空列で $q,S$ をそのまま返す．cons では head の出力を tail の入力にし，
binding context も左から右へ延長する．lookup または規則中の部分関数が失敗した場合，
その規則は適用できない．

\subsection{具体化と一般化}

$\Inst_q$ と $\InstDual_q$ は supply-indexed に量化 binder を fresh variable へ置き換え，
次 supply も返す．expression context と pattern-function signature の量化 capability binder は
capability variable へだけ instantiate される．利用側の demand を根拠に producer capability を
constructor や product へ変えることはできない．一方，data／pattern constructor と primitive の
scheme は通常の binder-supported structural instance を許す．

$\Gen$ は signature と context に自由な variable を除いて量化する．
\[
\Gen(\Sighat,\Gamma,\tau)=
\forall(\fcv(\tau)\setminus(\fcv(\Sighat)\cup\fcv(\Gamma))).
\forall(\ftv(\tau)\setminus(\ftv(\Sighat)\cup\ftv(\Gamma))).\tau.
\]
scheme binder は局所的であり，後段の substitution が導入した同番号の自由 variable を
捕捉しない．$\mathtt{let}$ は value の synthesis 直後の $S\Gamma,S\tau$ を一般化する．
